\documentclass{article}
\usepackage[T1]{fontenc}
\usepackage{lmodern}
\usepackage{graphicx}
\usepackage{amsmath,amssymb}
\usepackage[a4paper,margin=2cm]{geometry}
\usepackage[absolute,overlay]{textpos}
\setlength{\TPHorizModule}{1mm}
\setlength{\TPVertModule}{1mm}
\usepackage[indonesian]{babel}
\usepackage{hyperref}
\setlength{\emergencystretch}{2em}
% unit: Halaman 44

\begin{document}

% === Halaman 44 (python/native) ===

44

2. Serangan aktif (active attack) 

• penyerang mengintervensi komunikasi dan ikut mempengaruhi sistem untuk 

keuntungan dirinya 

• penyerang dapat mengubah aliran pesan seperti: 


- menghapus  sebagian cipherteks,  


- mengubah cipherteks, 


- menyisipkan potongan cipherteks palsu, 


- me-replay pesan lama, 


- mengubah informasi yang tersimpan, dsb 

• Contoh: man-in-the-middle attack

\end{document}
